\chapter{Marked finite quantum dilogarithms}
\label{v4-arith:mf:chapter}
\chaptermark{Marked finite quantum dilogarithms}

Consider two simple objects \(\mathbf1,\rho\) with
\[
 \rho\otimes\rho\simeq\mathbf1\oplus\rho.
\]
An additive strong tensor functor to ordinary finite-dimensional vector
spaces would give a positive integer \(D\) satisfying \(D^2=1+D\).
No such integer exists. Nevertheless, the rule has a realization by
finite bimodules, whose tensor product is balanced over an algebra
with two idempotents. Its associativity maps retain information that
the dimension equation forgets.

A real quadratic unit produces a family of such examples. The input
is a finite function obtained from a meromorphic quantum
dilogarithm. Its values become coefficients of explicit fusion maps.
A normalized network recovers those values from the category.
We construct the category over the field of its normalized arithmetic
values, then construct its finite bimodule realization. Along the way,
a rank-one correction separates a finite operator identity from the
categorical pentagon.

\section{The arithmetic function}
\label{v4-arith:mf:arithmetic}

Write \(\mathrm e(z)=\exp(2\pi iz)\).
For \(|q|\ne1\), define the \(q\)-Pochhammer symbol by
\[
 (x;q)_\infty=
 \begin{cases}
  \displaystyle\prod_{j\ge0}(1-xq^j),&|q|<1,\\[2mm]
  \displaystyle\prod_{j\ge1}(1-xq^{-j})^{-1},&|q|>1.
 \end{cases}
\]
The first product converges locally uniformly in \(x\).
The second defines a meromorphic function of \(x\).
For
\(\gamma=\left(\begin{smallmatrix}a_\gamma&b_\gamma\\
c_\gamma&d_\gamma\end{smallmatrix}\right)\in SL_2(\mathbb Z)\)
and integers \(m,n\), put
\begin{equation}\label{v4-arith:mf:modular-quotient}
 \Phi_{\gamma,m,n}(z;\tau)
 =
 \frac{(q^m\mathrm e(z);q)_\infty}
 {(\widetilde q^{\,n}
       \mathrm e(z/(c_\gamma\tau+d_\gamma));\widetilde q)_\infty},
 \qquad
 q=\mathrm e(\tau),\quad
 \widetilde q=\mathrm e(\gamma\tau).
\end{equation}
Initially \(\tau\notin\mathbb R\), and the quotient is interpreted
meromorphically in \(z\).

The modular quantum dilogarithm has a meromorphic continuation to
\[
 (z,\tau)\in\mathbb C^2,\qquad
 c_\gamma\tau+d_\gamma\notin\mathbb R_{\le0}.
\]
We use the continuation of
Radchenko--Wheeler, equation~(1) and Section~2.2.\footnote{
D.~Radchenko and C.~Wheeler,
\emph{Real quadratic fields and finite quantum dilogarithms I}.
Version~2. \href{https://arxiv.org/abs/2609.21892v2}{arXiv:2609.21892v2}.

The product conventions are in Appendix~A.1.
The arithmetic identities are in Theorems~1--2 and Section~4.
The fusion reconstruction is Theorem~12 and Appendix~B.}

For \(c_\gamma>0\), fix the Dedekind eta multiplier
\[
 \mu_\gamma=
 \frac{\eta(\gamma w)}
      {\eta(w)\sqrt{c_\gamma w+d_\gamma}},
 \qquad w\in\mathbb H,
\]
using the holomorphic square root with positive value on the
positive real axis. Here
\(\eta(w)=\mathrm e(w/24)\prod_{r\ge1}(1-\mathrm e(rw))\).
The quotient defining \(\mu_\gamma\) is independent of \(w\)
and is a twenty-fourth root of unity.

\subsection{A quadratic lattice and its finite group}

Fix a real quadratic field \(K\subset\mathbb R\).
Write \(x'\) for the other real embedding of \(x\in K\).
Let \(I\subset K\) be a free abelian subgroup of rank two.
Choose a totally positive unit \(\varepsilon>1\) satisfying
\(\varepsilon I=I\). Its norm is one.

Define
\begin{equation}\label{v4-arith:mf:group}
 G_I=(\varepsilon-1)^{-1}I/I,\qquad
 N=|G_I|=\operatorname{Tr}_{K/\mathbb Q}(\varepsilon)-2.
\end{equation}
Indeed, multiplication by \(\varepsilon-1\) identifies this group
with \(I/(\varepsilon-1)I\).
The absolute determinant is
\(\lvert\operatorname{Nm}(\varepsilon-1)\rvert
=\operatorname{Tr}(\varepsilon)-2>0\).

An oriented basis \((\beta_1,\beta_2)\) of \(I\) satisfies
\(\beta_1\beta'_2-\beta'_1\beta_2>0\).
Choose one with \(\beta_2\) totally positive, and write
\[
 \varepsilon\beta_1=a_\gamma\beta_1+b_\gamma\beta_2,\qquad
 \varepsilon\beta_2=c_\gamma\beta_1+d_\gamma\beta_2.
\]
The resulting matrix \(\gamma\) lies in \(SL_2(\mathbb Z)\).
For \(\tau=\beta_1/\beta_2\), one has
\[
 \gamma\tau=\tau,\qquad c_\gamma\tau+d_\gamma=\varepsilon,
 \qquad
 c_\gamma=\frac{\varepsilon-\varepsilon^{-1}}{\tau-\tau'}>0.
\]
Thus \(\tau\) is the attractive real fixed point, inside the domain
of the continued function.

For \(x\in G_I\), choose a representative in \(K\).
There are unique integers \(u,v\) with
\[
 (\varepsilon^{-1}-1)x=u\beta_1+v\beta_2.
\]
For a nonzero class define
\begin{equation}\label{v4-arith:mf:E}
 E_I(x)=
 \mu_\gamma\Phi_{\gamma,u,0}(x/\beta_2;\tau),
 \qquad E_I(0)=\sqrt{\varepsilon}>0.
\end{equation}
The continuation in
\eqref{v4-arith:mf:modular-quotient} precedes this evaluation.
At the real fixed point the two separate products need not converge.
Their formally matching arguments therefore cannot be cancelled.

The representative and oriented-basis independence of
\eqref{v4-arith:mf:E} follows from the modular cocycle and its
index-shift identities. We use the precise lattice formulation and
change-of-basis theorem of the reciprocity paper.\footnote{
\emph{Stark units for real quadratic fields and reciprocity laws},
4 October 2026, Definition~1, Theorems~1--2, Proposition~1
and Appendix~B,
\url{https://www.danrad.net/papers/stark-reciprocity.pdf}.}
In particular, the prescribed meromorphic function selects the
arithmetic values. The scalar identities below do not alone specify
that selection.

\subsection{The Gaussian and Fourier normalization}

Let \(\omega_I:I\times I\to\mathbb Z\) be the alternating form with
\(\omega_I(\beta_1,\beta_2)=1\), extended
\(\mathbb Q\)-bilinearly to \(K\).
For \(t\in I\), set
\[
 \vartheta_I(t)=
 \begin{cases}
  1,&t\in2I,\\
  -1,&t\notin2I.
 \end{cases}
\]
The Gaussian and its polarization are
\begin{align}
 Q_I(x)
 &=\vartheta_I((\varepsilon-1)x)
       \mathrm e\!\left(\tfrac12\omega_I(x,\varepsilon x)\right),
       \label{v4-arith:mf:gaussian}\\
 \chi_I(x,y)
 &=\frac{Q_I(x+y)}{Q_I(x)Q_I(y)}
   =\mathrm e\bigl(\omega_I(x-\varepsilon x,y)\bigr).
       \label{v4-arith:mf:chi}
\end{align}
Both are well defined on \(G_I\).

To see the descent of \(Q_I\), use
\[
 \vartheta_I(t+s)\vartheta_I(t)\vartheta_I(s)
       =(-1)^{\omega_I(t,s)},\qquad t,s\in I,
\]
and \(\omega_I(\varepsilon x,\varepsilon y)=\omega_I(x,y)\).
Replacing \(x\) by \(x+t\) changes the sign and exponential
by inverse factors. The same calculation gives
\eqref{v4-arith:mf:chi}.
One has \(Q_I(0)=1\), \(Q_I(-x)=Q_I(x)\), and all values
of \(Q_I\) are roots of unity.

The bicharacter \(\chi_I\) is nondegenerate.
Indeed,
\[
 \omega_I(x-\varepsilon x,y)
       =\omega_I(\varepsilon x,(\varepsilon-1)y).
\]
If its exponential is one for every
\(y\in(\varepsilon-1)^{-1}I\), unimodularity of \(\omega_I\)
gives \(\varepsilon x\in I\). Hence \(x\in I\).

For a function \(f:G_I\to\mathbb C\), use the Fourier transform
\[
 \widehat f(y)=\frac1{\sqrt N}
        \sum_{x\in G_I}f(x)\chi_I(x,-y),
\]
with the positive square root of \(N\).
Orthogonality of the characters gives
\(\widehat{\widehat f}(x)=f(-x)\).

\begin{theorem}[Arithmetic input]\label{v4-arith:mf:arithmetic-input}
Let \(E=E_I\), \(Q=Q_I\), and \(\chi=\chi_I\).
For a twenty-fourth root of unity \(\lambda\), the arithmetic
values satisfy
\begin{align}
 E(g)E(-g)&=Q(g)^{-1}\qquad(g\ne0),\label{v4-arith:mf:reflection}\\
 \widehat E(g)&=\lambda Q(g)E(g),\label{v4-arith:mf:fourier}\\
 \chi(x,y)E(x)E(y)
 &=-\widehat E(0)+\frac1{\sqrt N}
      \sum_t E(y-t)E(t)Q(t)E(x-t).
       \label{v4-arith:mf:cubic}
\end{align}
At zero the separate normalization is
\begin{equation}\label{v4-arith:mf:zero-normalization}
 E(0)^2=\sqrt N\,E(0)+1.
\end{equation}
Every \(E(g)\) is an algebraic unit.

Let \(m\) be the exponent of \(G_I\), and let
\(\mathcal O=(I:I)=\mathbb Z+c\mathcal O_K\).
The normalized values
\(e_{I,g}=E_I(g)/E_I(0)\) belong to
\(\mathcal O_L^\times\), where \(L=H^+_{\mathcal O,m}\)
is the narrow ray class field modulo \(m\mathcal O\).
For an integral invertible \(\mathcal O\)-ideal \(\mathfrak a\)
prime to \(cm\),
\begin{equation}\label{v4-arith:mf:artin-input}
 \operatorname{Art}(\mathfrak a)(e_{I,g})
       =e_{\mathfrak a^{-1}I,g}.
\end{equation}
The right-hand label is the class of the same representative.
The Artin symbol uses arithmetic Frobenius.
\end{theorem}

The scalar identities and algebraicity are the
Radchenko--Wheeler theorems.
The unit and reciprocity assertions are Theorems~1--2 of the
reciprocity paper. These are imported arithmetic results.
The identity \eqref{v4-arith:mf:zero-normalization} also follows
directly from
\(\varepsilon+\varepsilon^{-1}=N+2\).
It remains part of the input when \(N\le4\), where a bare cubic
equation can have additional solutions.

\section{The arithmetic coefficient field}
\label{v4-arith:mf:coefficients}

Suppress \(I\) in the notation. Put
\(u=\sqrt\varepsilon>0\) and \(s=\sqrt N>0\).
Then \(s=u-u^{-1}\). Define
\begin{equation}\label{v4-arith:mf:b}
 b_g=\frac{E(g)}s,\qquad
 b_0=\frac{\varepsilon}{\varepsilon-1}
      =\frac{\varepsilon-1}{N},\qquad b_g=b_0e_g.
\end{equation}

Consequently the coefficient field is
\[
 F_I=K(b_g:g\in G_I)=K(e_g:g\in G_I)\subseteq L.
\]
Neither \(u\) nor \(s\) needs to belong to \(F_I\).
Every coefficient used in the algebraic construction below does.

Set
\begin{equation}\label{v4-arith:mf:parameters}
 d=-b_0^{-1}=\varepsilon^{-1}-1,\qquad
 \beta=\sum_gb_g,\qquad
 a=\frac{\beta}{Nb_0},\qquad c_1=\frac{b_0}{\beta}.
\end{equation}
Fourier evaluation at zero gives \(\beta=\lambda u\ne0\).
Thus \(a=\lambda/s\), \(c_1=\lambda^{-1}/s\), and
\(ac_1=1/N\). The notation \(c_1\) distinguishes this coefficient
from the conductor \(c\).

Reflection recovers the Gaussian over the same field:
\[
 Q(0)=1,\qquad
 Q(g)=\frac1{Nb_gb_{-g}}\quad(g\ne0).
\]
Its polarization therefore has values in \(F_I\).
The arithmetic identities become
\begin{align}
 b_0^2-b_0&=1/N,\qquad d^2=N(d+1),\label{v4-arith:mf:scalar}\\
 \sum_x b_x\chi(x,-y)&=NaQ(y)b_y,\label{v4-arith:mf:fourier-b}\\
 \sum_t b_{y-t}b_tQ(t)b_{x-t}-\chi(x,y)b_xb_y
       &=\beta/N.\label{v4-arith:mf:cubic-b}
\end{align}
These follow by substituting \(E(g)=s b_g\).
In particular, the constant in
\eqref{v4-arith:mf:cubic-b} retains the Fourier normalization.

\begin{example}\label{v4-arith:mf:fibonacci}
Take \(I=\mathbb Z[\varepsilon]\) with
\(\varepsilon=(3+\sqrt5)/2\).
Then \(N=1\), \(G_I\) is trivial, and
\[
 b_0=\frac{1+\sqrt5}{2},\qquad
 d=\frac{1-\sqrt5}{2},\qquad a=c_1=1.
\]
The fusion rule will be \(\rho^2\simeq\mathbf1\oplus\rho\).
The roots \(b_0,d\) will give its positive dimension and the
specified duality closure, respectively.
\end{example}

\section{The Leavitt presentation}
\label{v4-arith:mf:leavitt}

Let \(\mathcal B_I\) be the unital \(F_I\)-algebra generated by
\[
 X_g,\ X'_g,\ T_g,\ T'_g,\qquad g\in G_I,
\]
with relations
\begin{align}
 X'_gX_h=T'_gT_h&=\delta_{g,h}1,&
 X'_gT_h=T'_gX_h&=0,\label{v4-arith:mf:leavitt-relations}\\
 \sum_g(X_gX'_g+T_gT'_g)&=1.&&\notag
\end{align}
Primes denote algebraic dual generators, without any
complex-conjugate or adjoint operation.
This is the Leavitt algebra on \(2N\) generator pairs.

For \(g\in G_I\), define
\begin{align*}
 \alpha_g(X_h)&=X_{g+h},&
 \alpha_g(X'_h)&=X'_{g+h},\\
 \alpha_g(T_h)&=\chi(g,h)^{-1}T_h,&
 \alpha_g(T'_h)&=\chi(g,h)T'_h.
\end{align*}
These formulas preserve the relations and give
\(\alpha_g\alpha_h=\alpha_{g+h}\).
Also define
\begin{equation}\label{v4-arith:mf:U}
 U_g=\sum_h\chi(g,h)^{-1}X_hX'_h
           +\sum_h T_{h-g}T'_h.
\end{equation}
Direct multiplication gives
\[
 U_0=1,\qquad U_gU_h=U_{g+h},\qquad
 \alpha_g(U_h)=\chi(g,h)U_h.
\]
Thus each \(U_g\) is invertible, with inverse \(U_{-g}\).

The noninvertible generator of the category is an endomorphism
\(\rho\) of \(\mathcal B_I\).
Its four families of images are
\begin{align}
 \rho(X_g)
 &=U_g\left(d^{-1}\sum_hX_h+
              \sum_h Q(h)T_hT_{-h}\right)U_{-g},
       \label{v4-arith:mf:rho-X}\\
 \rho(X'_g)
 &=d^{-1}U_g\left(\sum_hX'_h+
              \sum_h Q(h)^{-1}T'_{-h}T'_h\right)U_{-g},
       \label{v4-arith:mf:rho-Xprime}\\
 \rho(T_g)
 &=\frac ad\sum_{h,k}\chi(h-g,k)X_hT'_k
   +\frac{c_1}{Q(g)}
       \sum_{h,k}\chi(g+h,k)^{-1}T_kX_hX'_h
       \notag\\
 &\quad+\sum_{h,k}Q(h)b_{h+g}\chi(g,k)^{-1}
                       T_{h+k}T_{-h}T'_k,
       \label{v4-arith:mf:rho-T}\\
 \rho(T'_g)
 &=c_1\sum_{h,k}\chi(h-g,k)^{-1}T_kX'_h
   +aQ(g)\sum_{h,k}\chi(g+h,k)X_hX'_hT'_k
       \notag\\
 &\quad+\sum_{h,k}\frac{Q(g+h)}{Q(h)}
          b_{-h-g}\chi(g,k)T_kT'_{-h}T'_{h+k}.
       \label{v4-arith:mf:rho-Tprime}
\end{align}
All sums run over \(G_I\).
All four formulas are over \(F_I\).

\subsection{The tensor category and its duality}

For unital endomorphisms \(\sigma,\tau\) of \(\mathcal B_I\), put
\[
 \operatorname{Hom}(\sigma,\tau)=
 \{x\in\mathcal B_I:
       x\sigma(z)=\tau(z)x\text{ for every }z\in\mathcal B_I\}.
\]
Composition is multiplication, and tensor product of objects is
composition of endomorphisms. For
\(x:\sigma\to\sigma'\) and \(y:\tau\to\tau'\), set
\begin{equation}\label{v4-arith:mf:tensor-map}
 x\otimes y=x\sigma(y)=\sigma'(y)x:
                \sigma\tau\longrightarrow\sigma'\tau'.
\end{equation}
This gives a strict tensor category.
Let \(\mathcal C_I\) be the additive idempotent completion of its
subcategory generated by the \(\alpha_g\) and \(\rho\).

A fusion category is a rigid semisimple linear tensor category
with finitely many simple objects, finite-dimensional Hom spaces
and a simple tensor unit. It is split over \(F_I\) when the
endomorphism field of every simple object is \(F_I\).

\begin{theorem}\label{v4-arith:mf:fusion}
Equations \eqref{v4-arith:mf:rho-X}--\eqref{v4-arith:mf:rho-Tprime}
define a unital endomorphism.
The category \(\mathcal C_I\) is a split rigid fusion category
over \(F_I\).
Its simple objects are the pairwise distinct
\(\alpha_g\), \(g\in G_I\), and \(\rho\), with rules
\begin{equation}\label{v4-arith:mf:fusion-rule}
 \alpha_g\alpha_h=\alpha_{g+h},\qquad
 \alpha_g\rho\simeq\rho\simeq\rho\alpha_g,\qquad
 \rho^2\simeq\bigoplus_g\alpha_g\oplus\rho^{\oplus N}.
\end{equation}
For \(\rho^\vee=\rho\), one may choose coevaluation \(X_0\)
and evaluation \(dX'_0\).
\end{theorem}

\begin{proof}
Start with the complex Leavitt reconstruction of
Radchenko--Wheeler, Section~5.4, Theorem~12 and Appendix~B.
In its notation make the change
\[
 X_g=\sqrt d\,S_g,\qquad X'_g=S'_g/\sqrt d.
\]
The Leavitt relations are unchanged.
Substitution gives exactly
\eqref{v4-arith:mf:rho-X}--\eqref{v4-arith:mf:rho-Tprime}.
The two mixed coefficients in \(\rho(T_g)\) become \(a/d,c_1\).
Those in \(\rho(T'_g)\) become \(c_1,a\).
The quadratic terms in \(\rho(X_g)\) lose their square root,
and both terms in \(\rho(X'_g)\) acquire \(d^{-1}\).

The reconstruction theorem proves the defining relations and
\begin{equation}\label{v4-arith:mf:endo-identities}
 \alpha_g\rho=\rho,\qquad
 \rho\alpha_g=\operatorname{Ad}(U_g)\rho,\qquad
 \rho^2(z)=\sum_gX_g\alpha_g(z)X'_g+
                     \sum_hT_h\rho(z)T'_h.
\end{equation}
Faithful extension of scalars from \(F_I\) to \(\mathbb C\)
descends these equalities.

The same extension descends the Hom calculation.
The intertwiner equations need only be imposed on the \(4N\)
algebra generators. Each Hom space is consequently the kernel
of one \(F_I\)-linear map
\[
 \mathcal B_I\longrightarrow\mathcal B_I^{\,4N}.
\]
A field extension preserves this kernel, even though
\(\mathcal B_I\) is infinite-dimensional.
The complex calculation therefore gives
\(\operatorname{End}(\alpha_g)=\operatorname{End}(\rho)=F_I\)
and zero Hom spaces between distinct displayed simples.

The last formula in \eqref{v4-arith:mf:endo-identities}
splits \(\rho^2\), with inclusions \(X_g,T_h\)
and projections \(X'_g,T'_h\).
The Leavitt relations give their orthogonality and completeness.
Every tensor word therefore splits into the listed simples.
Finite sums and idempotent summands give a split semisimple
category with finite-dimensional Hom spaces.

The dual of \(\alpha_g\) is \(\alpha_{-g}\), with both maps \(1\).
For \(\rho\), the two snake equations are
\[
 dX'_0\rho(X_0)=1,\qquad \rho(dX'_0)X_0=1.
\]
For the first, the \(X_0\) term contributes
\(d\,d^{-1}=1\), and every \(T\)-term vanishes.
The second follows by the dual contraction in
\eqref{v4-arith:mf:rho-Xprime}.
This proves rigidity over the stated field.
\end{proof}

For the fusion rules above, the Frobenius--Perron dimension is
the positive dimension function assigning dimension one to each
invertible simple object.

\begin{proposition}\label{v4-arith:mf:dimension}
The Frobenius--Perron dimension of \(\rho\) is
\(\delta=\varepsilon-1=Nb_0\), while the specified duality
closure of \(1_\rho\) is \(d=\varepsilon^{-1}-1<0\).
Moreover,
\[
 \varepsilon=1+Nb_0=\frac{b_0}{b_0-1},\qquad
 K=\mathbb Q(b_0).
\]
There is no additive strong tensor functor from \(\mathcal C_I\)
to ordinary finite-dimensional vector spaces over any field.
\end{proposition}

\begin{proof}
The positive dimension satisfies \(\delta^2=N(1+\delta)\).
The identity \((\varepsilon-1)^2=N\varepsilon\)
selects \(\delta=\varepsilon-1\).
The other root is \(d\).
The chosen cup and cap have composite \(dX'_0X_0=d\).
Formula \eqref{v4-arith:mf:b} recovers the period and its
quadratic field.

Under the proposed tensor functor, every invertible \(\alpha_g\)
would have dimension one.
The image of \(\rho\) would have a nonzero integer dimension \(D\),
since its square contains the tensor unit.
Then \(D^2=N(D+1)\).
It follows that \(D+1\) divides \(D^2\), hence divides \(1\),
which is impossible for \(D>0\).
No faithfulness assumption is needed.
\end{proof}


\subsection{An integral coefficient algebra}

Put \(\mathcal S=\mathcal O_L[1/N]\).
The normalized units \(e_g\) and
\(b_0(b_0-1)=1/N\) show that every \(b_g\) is a unit of
\(\mathcal S\). Also \(\beta=\lambda u\) is an algebraic unit.
Since \(\beta\in L\), it belongs to \(\mathcal O_L^\times\).
Thus \(d,a,c_1,Q(g)\) and all their required inverses lie in
\(\mathcal S\).

\begin{proposition}\label{v4-arith:mf:integral}
The Leavitt presentation and all four formulas for \(\rho\)
are defined over \(\mathcal S\).
The underlying Leavitt algebra is free as an \(\mathcal S\)-module.
Its extension to \(L\) agrees with the extension of the
\(F_I\)-presentation to \(L\).
\end{proposition}

\begin{proof}
For the freeness assertion, combine the \(X\) and \(T\) generators
into \(s_1,\ldots,s_{2N}\), with dual generators
\(t_1,\ldots,t_{2N}\).
Use the reduction rules
\[
 t_i s_j\longmapsto\delta_{ij},\qquad
 s_{2N}t_{2N}\longmapsto
       1-\sum_{i<2N}s_it_i.
\]
Choose a length and lexicographic order with \(s_{2N}\)
larger than the other positive generators.
Every reduction decreases this order.
The irreducible words are \(1\), pure positive or negative words,
and positive words followed by negative words without the
forbidden middle pair \(s_{2N}t_{2N}\).

The only nontrivial overlap types are
\(t_i s_{2N}t_{2N}\) and \(s_{2N}t_{2N}s_j\).
Reducing either end of the first gives
\(\delta_{i,2N}t_{2N}\).
Reducing either end of the second gives
\(\delta_{2N,j}s_{2N}\).
Disjoint reductions commute. Hence the irreducible words give a
basis over any coefficient ring.

All displayed coefficients lie in \(\mathcal S\).
The free normal form embeds this algebra into its scalar extension
over \(L\). The endomorphism identities already hold there,
so they hold over \(\mathcal S\).
\end{proof}

This integral model permits specialization away from \(N\).
It makes no semisimplicity assertion about the resulting category
over a residue field. Primes dividing \(N\) require an additional
integral construction.

\section{Fusion bases and the arithmetic network}
\label{v4-arith:mf:network-section}

The fusion rule alone does not specify associativity coefficients.
Here the Leavitt intertwiners give complete labeled splitting maps.
For simple objects \(i,j,k\), let
\(v_{ij}^{k,\mu}:k\to i\otimes j\) and
\(\bar v_{ij}^{k,\mu}:i\otimes j\to k\)
denote dual inclusions and projections.
Choose them as follows:
\[
\begin{array}{c|c|c}
 i\otimes j,\ k&v_{ij}^{k,\mu}&\bar v_{ij}^{k,\mu}\\ \hline
 \alpha_g\alpha_h,\ \alpha_{g+h}&1&1\\
 \alpha_g\rho,\ \rho&1&1\\
 \rho\alpha_g,\ \rho&U_g&U_{-g}\\
 \rho^2,\ \alpha_g&X_g&X'_g\\
 \rho^2,\ \rho\text{ in channel }h&T_h&T'_h
\end{array}
\]
They satisfy
\[
 \bar v_{ij}^{k,\mu}v_{ij}^{l,\nu}
       =\delta_{kl}\delta_{\mu\nu}1_k,\qquad
 \sum_{k,\mu}v_{ij}^{k,\mu}\bar v_{ij}^{k,\mu}
       =1_{i\otimes j}.
\]

For four simple objects \(i,j,k,l\), form bases of
\(\operatorname{Hom}(l,i\otimes j\otimes k)\):
\[
 L_{p,\mu,\nu}
   =(v_{ij}^{p,\mu}\otimes1_k)v_{pk}^{l,\nu},\qquad
 R_{q,\alpha,\beta}
   =(1_i\otimes v_{jk}^{q,\alpha})v_{iq}^{l,\beta}.
\]
Their dual projections are denoted by \(\bar L,\bar R\).
The full recoupling matrix is
\begin{equation}\label{v4-arith:mf:F}
 F^l_{ijk}[(q,\alpha,\beta),(p,\mu,\nu)]
       =\bar R_{q,\alpha,\beta}L_{p,\mu,\nu},
 \qquad
 (F^{-1})[(p,\mu,\nu),(q,\alpha,\beta)]=\bar L_{p,\mu,\nu}R_{q,\alpha,\beta}.
\end{equation}
Each entry is a scalar in \(\operatorname{End}(l)=F_I\).
Formula \eqref{v4-arith:mf:tensor-map} and Leavitt reduction
compute every entry from the stated generators.

The completeness identities give \(F^{-1}F=FF^{-1}=1\).
For four tensor factors, every composite of successive changes
of basis equals the direct change between its first and last
fusion trees. The two routes around the pentagon therefore agree.
The unit choices give the triangle equations in the same way.
A strict endomorphism presentation thus produces nontrivial
recoupling matrices on its chosen simple-object skeleton.

\subsection{Recovering the normalized values}

Consider the endomorphism of \(\rho\)
\begin{equation}\label{v4-arith:mf:network}
 \mathcal N_g
 =T'_0(T'_0\otimes1_\rho)(1_\rho\otimes T_g)T_0
 =T'_0T'_0\rho(T_g)T_0.
\end{equation}
The second expression is in the strict presentation.
In a general presentation the first expression includes its
specified associativity maps.

\begin{proposition}\label{v4-arith:mf:network-value}
The network satisfies
\[
 \mathcal N_g=b_g1_\rho,\qquad
 \mathcal N_g\mathcal N_0^{-1}=e_g1_\rho.
\]
The scalars are intrinsic to the category with labeled invertible
simples. They are preserved by linear tensor equivalences respecting
those labels.
\end{proposition}

\begin{proof}
Substitute \eqref{v4-arith:mf:rho-T}.
Every mixed \(X\)-term vanishes.
The right contraction forces \(k=0\), and the next forces \(h=0\).
The remaining coefficient is \(b_g\).

To recover the marking, choose isomorphisms
\(a_g:\alpha_g\otimes\rho\to\rho\).
The equations
\[
 (a_g\otimes1_\rho)(1_{\alpha_g}\otimes T_0)=T_0a_g
       \quad\text{for every }g
\]
define a line in \(\operatorname{Hom}(\rho,\rho^2)\).
In the explicit basis the action is
\(T_h\mapsto\chi(g,h)^{-1}T_h\).
Nondegeneracy of \(\chi\) makes the common invariant space
exactly \(F_IT_0\).

Choose \(u_g:\rho\to\rho\otimes\alpha_g\).
Normalize it so that
\[
 (1_\rho\otimes a_g)(u_g\otimes1_\rho)
\]
acts identically on \(\operatorname{Hom}(\mathbf1,\rho^2)\).
This fixes the endomorphism \(U_g\).

Set \(T_g=U_{-g}T_0\), and choose the dual projections \(T'_g\).
Only a common rescaling of the \(T_g\) remains.
Its two factors cancel the two inverse factors in
\eqref{v4-arith:mf:network}.
The construction is therefore preserved by the stated equivalences.
\end{proof}

The unit summand normalizes \(U_g\), but \(T_0\) itself belongs to
\(\operatorname{Hom}(\rho,\rho^2)\).
These are different channels. The specified duality closure is
\begin{equation}\label{v4-arith:mf:closure}
 dX'_0(1_\rho\otimes\mathcal N_g)X_0
       =d\,b_g=-e_g.
\end{equation}
Thus the closure recovers a normalized arithmetic unit, with its
negative sign.

This closure cannot be replaced by a normalized cyclic functional
on the ambient Leavitt algebra.
Indeed,
\[
 \sum_{i=1}^{2N}[t_i,s_i]=(2N-1)1.
\]
In characteristic zero, the unit has zero image in
\[
 HH_0(\mathcal B_I)=
       \mathcal B_I/[\mathcal B_I,\mathcal B_I].
\]
Every linear functional vanishing on commutators therefore sends
\(1\) to zero. The marked categorical closure in
\eqref{v4-arith:mf:closure} has a different domain.

Relabelling invertible simples relabels \(b_g\).
A semilinear tensor equivalence applies its field embedding to each
network coefficient. Recovering the unnormalized \(E(g)\) also
requires the chosen \(\sqrt N\). The normalized coefficients
and their ratios are already defined over \(F_I\).

\section{The finite Heisenberg operators}
\label{v4-arith:mf:operators}

The multiplicity space
\[
 H_I=\operatorname{Hom}(\rho,\rho^2)
       =\bigoplus_{h\in G_I}F_IT_h
\]
carries two group actions:
\begin{equation}\label{v4-arith:mf:heisenberg}
 p^g(T_h)=U_gT_h=T_{h-g},\qquad
 q^g(T_h)=\alpha_{-g}(T_h)=\chi(g,h)T_h.
\end{equation}
Their commutation rule is
\[
 p^gq^h=\chi(g,h)q^hp^g.
\]
Because \(Q(g+h)=Q(g)Q(h)\chi(g,h)\), the formula
\[
 (p+q)^g=Q(g)q^gp^g
\]
also gives a group action. Here \(p+q\) denotes this action,
rather than the sum of two operators.

For any group action \(r\) and function \(f:G_I\to F_I\), define
\begin{equation}\label{v4-arith:mf:functional-calculus}
 f(r)=\frac1N\sum_{x,y\in G_I}f(y)\chi(x,-y)r^x.
\end{equation}
Character orthogonality gives \(f(q)T_h=f(h)T_h\).
Let \(\delta_0\) denote the indicator of zero, and put
\[
 P_0=\delta_0(q),\qquad
 P_{\mathrm{av}}=\delta_0(p),\qquad D=P_0P_{\mathrm{av}}.
\]
The first two are projections, respectively onto \(F_IT_0\)
and the invariant vectors for translation.
Their product acts by
\begin{equation}\label{v4-arith:mf:D}
 D\!\left(\sum_hv_hT_h\right)
       =\frac1N\left(\sum_hv_h\right)T_0,\qquad
 D^2=\frac DN,\qquad \operatorname{Tr}(D)=\frac1N.
\end{equation}
For \(N>1\), the rank-one operator \(D\) is not an idempotent.

Although \(E\) need not be \(F_I\)-valued, its Fourier transform is:
\[
 \eta(g):=\widehat E(g)=NaQ(g)b_g\in F_I.
\]

\begin{theorem}\label{v4-arith:mf:corrected-identity}
On \(H_I\), one has
\begin{equation}\label{v4-arith:mf:corrected-formula}
 \eta(p)\eta(q)-\eta(q)\eta(p+q)\eta(p)=-N\beta D.
\end{equation}
The tensor-unit projection of \(\rho^2\) acts as zero on \(H_I\),
whereas \(D\ne0\).
\end{theorem}

\begin{proof}
Fourier inversion and \eqref{v4-arith:mf:functional-calculus} give
\[
 \eta(p)=\sum_xb_xp^{-x},\qquad
 \eta(q)=\sum_yb_yq^{-y},\qquad
 \eta(p+q)=\sum_tb_tQ(t)q^{-t}p^{-t}.
\]
Order every term as \(q^{-y}p^{-x}\).
The coefficient of that term in the difference is
\[
 \chi(x,y)b_xb_y
       -\sum_tb_{y-t}b_tQ(t)b_{x-t}
       =-\frac{\beta}{N},
\]
by \eqref{v4-arith:mf:cubic-b}.
On the other hand,
\[
 D=\frac1{N^2}\sum_{x,y}q^{-y}p^{-x}.
\]
This proves the factor \(N\) in
\eqref{v4-arith:mf:corrected-formula}.

The tensor-unit projection is \(X_0X'_0\).
For every \(v:\rho\to\rho^2\),
the composite \(X'_0v\) lies in
\(\operatorname{Hom}(\rho,\mathbf1)=0\).
Its action on \(H_I\) vanishes.
By contrast, \(D(T_0)=T_0/N\ne0\).
\end{proof}

\begin{example}\label{v4-arith:mf:order-two}
Let \(I=\mathbb Z[\sqrt3]\), \(\varepsilon=2+\sqrt3\),
and orient \(I\) by \((\sqrt3,1)\).
The nonzero class of \(G_I\simeq\mathbb Z/2\)
is represented by \(x=(\sqrt3-1)/2\).
Since
\[
 (\varepsilon-1)x=1,\qquad
 \omega_I(x,\varepsilon x)=\frac12,
\]
formula \eqref{v4-arith:mf:gaussian} gives \(Q(x)=-i\).

Reflection gives \(2b_1^2=i\).
Among its two signs, the Fourier condition
\(\lvert\beta\rvert=\sqrt\varepsilon\) selects
\[
 b_0=\frac{1+\sqrt3}{2},\qquad
 b_1=\frac{1+i}{2},\qquad
 \beta=1+\frac{\sqrt3+i}{2}.
\]

In the basis \(T_0,T_1\),
\[
 \eta(p)=
 \begin{pmatrix}b_0&b_1\\b_1&b_0\end{pmatrix},\qquad
 \eta(q)=
 \begin{pmatrix}\beta&0\\0&b_0-b_1\end{pmatrix}.
\]
Put \(Q_q=\operatorname{diag}(1,-i)\).
Then \(\eta(p+q)=Q_q^{-1}\eta(p)Q_q\), and multiplication yields
\[
 \eta(p)\eta(q)-\eta(q)\eta(p+q)\eta(p)
  =-\beta\begin{pmatrix}1&1\\0&0\end{pmatrix}
  =-2\beta D.
\]
In particular, both the chosen multiplicity space and the
normalization of the correction are visible in two dimensions.
\end{example}

The identity \eqref{v4-arith:mf:corrected-formula}
is not an associativity failure of \(\mathcal C_I\).
Its recoupling matrices satisfy the categorical pentagon.
Moreover, the nonzero correction cannot be the image of a boundary
under a chain map to \(\operatorname{End}(H_I[0])\).
That target has no term in degree \(-1\).
A comparison with supported chiral operations must retain a larger
complex and specify its reduction to this multiplicity space.

\section{Finite bimodule realization}
\label{v4-arith:mf:finite}

We now construct the finite realization promised at the beginning.
Let \(k\) be a field of characteristic zero, and let
\(\mathcal C\) be a split fusion category over \(k\).
Choose simple representatives \((S_i)_{i\in\mathcal I}\),
with \(S_0=\mathbf1\), and set
\[
 A=\bigoplus_{i\in\mathcal I}ke_i,\qquad
 \omega(X)_{ij}=
       \operatorname{Hom}_{\mathcal C}(S_i,X\otimes S_j).
\]
Declare \(e_i\omega(X)e_j=\omega(X)_{ij}\).
For \(f:X\to X'\), the map \(\omega(f)\) is postcomposition
with \(f\otimes1\).
Each \(\omega(X)\) is a finite-dimensional \(A\)-bimodule.

Write
\[
 a_{X,Y,Z}:(X\otimes Y)\otimes Z\longrightarrow
                   X\otimes(Y\otimes Z)
\]
for the associator of \(\mathcal C\).
For \(u:S_i\to X\otimes S_j\) and
\(v:S_j\to Y\otimes S_k\), define
\begin{equation}\label{v4-arith:mf:tensorator}
 \Phi_{X,Y}(u\otimes_Av)
       =a^{-1}_{X,Y,S_k}(1_X\otimes v)u.
\end{equation}
The output belongs to
\(\operatorname{Hom}(S_i,(X\otimes Y)\otimes S_k)\).
The balanced tensor retains precisely the matching middle labels.

\begin{theorem}\label{v4-arith:mf:finite-functor}
The maps \(\Phi_{X,Y}\) are isomorphisms.
Together with \(\omega(\mathbf1)\simeq A\), they make
\[
 \omega:\mathcal C\longrightarrow\operatorname{Bimod}^{\mathrm{fd}}_A
\]
a faithful exact strong monoidal functor.
It need not be full.
\end{theorem}

\begin{proof}
For each \(Y\otimes S_k\), choose splitting maps
\[
 v_{j\alpha}:S_j\longrightarrow Y\otimes S_k,\qquad
 w_{j\alpha}:Y\otimes S_k\longrightarrow S_j
\]
with
\[
 w_{j\alpha}v_{l\beta}
       =\delta_{jl}\delta_{\alpha\beta}1_{S_j},
 \qquad \sum_{j,\alpha}v_{j\alpha}w_{j\alpha}=1.
\]
The inverse to \eqref{v4-arith:mf:tensorator} is
\begin{equation}\label{v4-arith:mf:tensorator-inverse}
 h\longmapsto
 \sum_{j,\alpha}
 (1_X\otimes w_{j\alpha})a_{X,Y,S_k}h\otimes_Av_{j\alpha}.
\end{equation}
Composing in one order uses the sum of the splitting idempotents.
Composing in the other uses their pairwise orthogonality.
A change of splitting basis acts by inverse matrices on the two
factors, so the displayed sum is independent of that choice.

For three objects, both iterated tensorators compose the same
three maps from simple objects.
The associator pentagon equates the two parenthesizations.
The same argument handles every finite tensor string, and the
unitors give \(\omega(\mathbf1)\simeq A\).
The column \(j=0\) is
\(\bigoplus_i\operatorname{Hom}(S_i,X)\).
It detects every morphism out of \(X\), proving faithfulness.
Split semisimplicity gives exactness.

Finally, \(\operatorname{End}_{A-A}(A)=A\),
whereas \(\operatorname{End}_{\mathcal C}(\mathbf1)=k\).
For more than one simple label, the induced map on these endomorphism
spaces is not surjective.
\end{proof}

\subsection{An associator detected by its endpoint}

Take the pointed category with labels \(\mathbb Z/2\)
and scalar associator
\[
 c(g,h,l)=(-1)^{ghl}.
\]
Its pentagon is the mod-two three-cocycle equation.

The bimodule \(\omega(g)\) has one-dimensional entries at
\(i=g+j\). On the component with right endpoint \(r\),
the tensorator is multiplication by \(c(g,h,r)^{-1}\).
Its coherence equation is
\[
 c(g,h,l)c(g,h,l+r)^{-1}c(g+h,l,r)^{-1}
       =c(h,l,r)^{-1}c(g,h+l,r)^{-1}.
\]
At \(g=h=l=1\), forgetting the endpoint factors would leave
the false equality \(-1=1\).
The strict ambient category of bimodules thus retains the
nontrivial source associator through its tensorator.

\subsection{The finite arithmetic detector}

Apply the construction to \(\mathcal C_I\) over \(k=F_I\),
with labels \(\mathcal I=G_I\sqcup\{\rho\}\).
Let \(W_X\) be the matrix with entries
\((W_X)_{ij}=\dim_k\omega(X)_{ij}\).
The fusion rule gives
\begin{equation}\label{v4-arith:mf:multiplicities}
 W_{\alpha_g}=P_g\oplus[1],\qquad
 W_\rho=
 \begin{pmatrix}
  0_{N\times N}&\mathbf1\\
  \mathbf1^t&N
 \end{pmatrix},
\end{equation}
where \(P_g\) is the translation permutation on \(G_I\).
Matrix multiplication records the balanced tensor product:
\[
 W_\rho^2=\sum_gW_{\alpha_g}+NW_\rho,\qquad
 \dim_k\omega(\rho)=3N.
\]

For \(N=2\), this is
\[
 W_\rho=\begin{pmatrix}0&0&1\\0&0&1\\1&1&2\end{pmatrix}.
\]
Its entries sum to \(6\), while the entries of \(W_\rho^2\)
sum to \(18\).
The ordinary vector-space tensor square would have dimension \(36\).

Transport every map in \eqref{v4-arith:mf:network}
through \(\omega\), using \(\Phi\) for each tensor product.
The resulting endomorphism \(M_g\) of \(\omega(\rho)\) satisfies
\begin{equation}\label{v4-arith:mf:matrix-detector}
 M_g=b_g1_{\omega(\rho)},\qquad
 b_g=\frac{\operatorname{Tr}_k(M_g)}{3N},\qquad
 E(g)=\frac{\sqrt N}{3N}\operatorname{Tr}_k(M_g).
\end{equation}
The first recovery takes place over \(F_I\).
The second uses the chosen radical. Ratios recover \(e_g\)
without it.

This matrix trace works because the specified composite ends in
the simple object \(\rho\) and is scalar there.
It does not identify ordinary matrix trace with the duality
closure \eqref{v4-arith:mf:closure} on arbitrary diagrams.
The finite realization also keeps its chosen categorical
intertwiners, rather than every bimodule map.

Finally, it is not a finite-dimensional unital representation of
\(\mathcal B_I\).
Such a representation on \(V\) would give
\(V\simeq V^{\oplus2N}\) from the Leavitt generators.
Equivalently, matrix trace applied to
\(\sum_i[t_i,s_i]=(2N-1)1\) gives
\(0=(2N-1)\dim_kV\), impossible in characteristic zero.

The arithmetic values have therefore been retained by actual
fusion maps, a normalized network and a finite balanced tensor.
A further field realization must transport these maps and their
duality, rather than only the scalar detector.
Its supported operation complex must also explain how the
correction \(-N\beta D\) arises after reduction.
